\documentclass[10pt,a4paper]{amsart}
\usepackage[margin=19mm]{geometry}
\usepackage{amsmath,amssymb,mathtools}
\usepackage[scr=rsfs,cal=euler]{mathalfa}
\usepackage{xfrac}
\usepackage[unicode,hidelinks,bookmarks=false]{hyperref}
\newcommand{\ie}{i.e.\ }
\newcommand{\cf}{cf.\ }
\newcommand{\resp}{respectively\ }
\newcommand{\Subsection}{\subsection}
\providecommand{\nobreakdash}{\nobreak\hbox{-}}
\setlength{\emergencystretch}{2em}
\multlinegap0pt
\usepackage{amssymb}
\newtheorem{theorem}{Theorem}
\title{On modulo-recurrence and window complexity in infinite words}
\author{Julien Cassaigne, Idrissa Kabor\'e}
\date{Source: arXiv:2606.13744v1 (CC BY 4.0)}
\begin{document}
\maketitle

\noindent\emph{Source and licence.} On modulo-recurrence and window complexity in infinite words. Version arXiv:2606.13744v1 (CC BY 4.0), distributed by arXiv under the Creative Commons Attribution 4.0 International licence, http://creativecommons.org/licenses/by/4.0/, as recorded in the arXiv licence metadata for this exact version and shown on the abstract page https://arxiv.org/abs/2606.13744v1. Full licence notice: \texttt{licenses/arxiv-2606.13744-CC-BY-4.0.txt}.\par\smallskip

\noindent\textit{Editorial scope.} Counted unit: the article's theorem liminf-fini=periodic (source0.tex lines 625--631). The article's complete original proof of it is delivered with it (source0.tex lines 633--680, one \texttt{proof} environment of 48 lines). No mathematics is written, repaired or re-derived: the statement and the proof are the article's own lines, in the article's own notation, and no retained line is substituted or re-flowed.  No further source-local context is needed: every cross-reference of every retained line resolves inside the two delivered spans.  Nothing was omitted from any retained span, so the package carries no omission record.  No retained line contains a citation, so the wrapper carries no bibliography and no bibliography entry.  No retained line contains a figure environment, an image-inclusion command or drawing code, so no image is delivered and the package declares no asset dependency.  Editorial formatting only, in addition: an AMS \texttt{amsart} preamble that restates the article's own declarations and macros used by the retained lines, namely the environments \texttt{theorem} on separate counters, together with the standard packages those lines need.  The retained environments print in this document as Theorem 1 in order of appearance, whereas the article prints the same items with the numbering of its own full document.  The complete native source of the article as distributed in the arXiv e-print for this exact version is bundled unchanged as \texttt{sources/202724/source0.tex}.
\begin{theorem} \label{liminf-fini=periodic}

Let $u$ be a uniformly recurrent word such that
 $$\liminf_{n\rightarrow \infty} P^{w}_{u}(n) P^{w}_{u}(n+1) <\infty/.$$
 Then,  $u$ is periodic.

\end{theorem}

\begin{proof}
Let $u$ be a uniformly recurrent word. If $\liminf P^{w}_{u}(n) P^{w}_{u}(n+1) <\infty$ then,  there exist an integer $M$ and an  increasing sequence $(n_k)$ such that $ P^{w}_{u}(n_{k}) P^{w}_{u}(n_{k}+1) \leq M$. %It follows $ P^{w}_{u}(n_{k})\leq M$ and $ P^{w}_{u}(n_{k}+1)\leq M$.

% In this stape, as in the uniformely recurrent case we show   that $ P^{w}_{u}(n_{k}) $ admits periods with arbitrary size; also to $ P^{w}_{u}(n_{k}+1)$. Thus $u$ is periodic.



%Let $u$ be a uniformly recurrent aperiodic word. Suppose that $P^w(n)$ is bounded, $ie$
%$$\exists M>0: \forall n\in \mathbb{N}, P^w(n)\leq M.$$
Let  $N$ be a term of   $(n_k)$ large enough. One has $P^w(N)\leq M$ and $P^w(N+1)\leq M$. Recall that  $P^w(n)=\#\mathcal{F}^{w}_{n}(u)$. For all $i\in \left\{0,1, \, \ldots, \, N-1\right\}$, let us define the words $x_i$, $y_i$, and $z_i$ as follows:
$$\begin{array}{lcl}

x_i&=&u_{Ni}\cdots u_{Ni+N-1}\\
y_i&=&u_{N\left(i+1\right)}\cdots u_{N\left(i+1\right)+N-1}\\
z_i&=&u_{\left(N+1\right)i}\cdots u_{\left(N+1\right)i+N}\\

\end{array}$$

Then, we have $ x_i,\, y_i\in \mathcal{F}^{w}_{N}\left(u\right)$, $ z_i\in \mathcal{F}^{w}_{N+1}\left(u\right) $ and $z_i$ is a factor of $x_iy_i$. Furthermore, observe that $z_i$ appears in $x_iy_i$ at position $i$.


Thus, it results the following inclusion
$$E=\left\lbrace\left(x_i,y_i, z_i\right):  0\leq i\leq N-1\right\rbrace \subseteq\mathcal{F}_{N}^{w}(u)\times \mathcal{F}_{N}^{w}(u)\times \mathcal{F}_{N+1}^{w}(u).$$


So, $\# E \leq M^3$. Therefore, there exists some triple $(x, y, z)\in E$  for which we have $\# I\geq \frac{N}{M^{3}}$ where $I=\left\lbrace i\in \left\lbrace 0,\, \cdots,\, N-1\right\rbrace : \left( x_i,y_i,z_i\right) =\left( x,y,z\right)  \right\rbrace $. Let us denote by  $|xy|_z$ the number of occurrences of $z$ in $xy$. Since,  $i\in I$  means that $z$ appears in $xy$ at position $i$, we have $|xy|_z\geq \frac{N}{M^3}$ and also
$$\min\left\{j-i: i,\, j\in I, \, i<j\right\}\leq \frac{N-1}{\#I-1}\leq \frac{N-1}{\frac{N}{M^3}-1}$$ 

Now, assume that $N>M^6$. Then $\frac{N}{M^3}-M^3>0$ and  $(\frac{N}{M^3}-1)(M^3+1)>N-1$. Thus, for $N>M^6$ we have $\min\left\{j-i: i,\, j\in I, \, i<j\right\}<M^3+1$.
 %and since $\left(\frac{N}{M^3}-1\right)\left(M^3+1\right)>N-1$ if $N> M^6$.

In other terms,  there exist $i,\,j\in I$ such that $0<j-i\leq M^3$~. Consequently, $z$ is $(j-i)$-periodic. Thus, we have

% So, $|xy|_z\geq \frac{N}{M^3}$ and
%
%$$\min\left\{j-i: i,\, j\in I, \, i<j\right\}\leq \frac{N-1}{\#I-1}\leq \frac{N-1}{\frac{N}{M^3}-1}<M^3+1$$
%since $\left(\frac{N}{M^3}-1\right)\left(M^3+1\right)>N-1$ if $N> M^6$.
%
%Then, there exist $i,\,j\in I$ such that $0<j-i\leq M^3$~. Consequently, $z$ is $(j-i)$-periodic. Thus, we have
$$\forall N>M^6,\exists q\in \mathcal{F}(u): 1\leq |q|\leq M^3\ \textrm{and}\  q^{\left\lfloor \frac{N+1}{\left|q\right|}\right\rfloor}\in \mathcal{F}(u).$$
As there are finitely many such $q$, then there exists $q\in \mathcal{F}(u),\, 1\leq |q|\leq M^3$ and an  infinite subset $K$ of $\mathbb{N}$ such that $q^k$ occurs in $u$ for all $k\in K$.
 
So
$$ \exists q\in L(u),\,   \forall k\in \mathbb{N},\ q^k \in L(u).$$


Thus, either $L(u)=L(q^\omega) $ and then $u$ is $|q|$-periodic; or there exists  $r\in L(u)-L(q^\omega)$ and $r\notin L(q^k)$ for all $k\in \mathbb{N}$. 
This latter  case contradicts uniform recurrence of $u$.  \end{proof}

\end{document}
